\ifthenelse{\equal{\AiAckOption}{true}}{%
    \ifthenelse{\equal{\LanguageOption}{portuguese}}{%
        \chapter*{Declaração sobre o Uso de Inteligência Artificial}
    
        \vspace{2em}
        
Reconheço a utilização integrada da inteligência artificial Gemini (\url{https://gemini.google.com}) em conjunto com o ecossistema e agente autónomo de desenvolvimento Google Antigravity (\url{https://antigravity.google}) como ferramentas de suporte no processo de programação, desenho arquitetural e automação de testes lógicos da aplicação \textit{Plot in a Pot}. A inteligência artificial operou como assistente no espaço de trabalho local seguro (\textit{sandbox}), executando tarefas de refatorização, criação e análise de dependências sob estrita validação e supervisão humana.

\vspace{1em}

\noindent\textbf{Descrição da Utilização}

\prompt[1]{Analisa a estrutura de dados atual do interpretador visual e implementa uma rotina modular capaz de ler o formato de ficheiros de texto do motor de jogo, isolando a lógica de transição dos nós para evitar quebras no processamento da árvore narrativa.}

\aioutput[1]{[O agente do Antigravity acedeu ao ambiente local através do ciclo de planeamento e execução, estruturando os módulos iniciais do interpretador de dados e aplicando funções de higienização de texto para garantir a integridade dos caminhos e ligações visuais.]}

\vspace{1em}

\prompt[2]{Executa a análise de impacto estrutural no código da aplicação caso o serviço de backend mude o motor de inferência local assente em microsserviços externos para uma execução nativa baseada em computação gráfica no lado do cliente.}

\aioutput[2]{[O ecossistema inspecionou as dependências do repositório, mapeou o fluxo de dados dos serviços de inteligência artificial e reestruturou os métodos para gerir o ciclo de vida do novo motor diretamente no navegador da Web, sugerindo melhorias na uniformização dos identificadores estruturais.]}

\vspace{1em}

\noindent\textbf{Processo de Validação Humana e Decisão Prática:}
A utilização do ecossistema Antigravity e do Gemini serviu para acelerar o desenvolvimento de protótipos e mapear o impacto de alterações estruturais na aplicação. No entanto, todas as decisões finais de engenharia de software, modelação definitiva da base de dados de localização e otimização dos algoritmos de validação do fluxo narrativo foram revistas, corrigidas e implementadas manualmente pelos autores. A arquitetura computacional final e a sua conformidade técnica permanecem sob o total controlo e responsabilidade intelectual da equipa de desenvolvimento humana.

    }{\ifthenelse{\equal{\LanguageOption}{spanish}}{%
        \chapter*{Declaración sobre el Uso de Inteligencia Artificial}

        \guideinfo{Es una buena práctica académica indicar brevemente cómo, por qué y cuándo se utilizó la IA. Explicar las acciones tomadas y las razones detrás de ellas fomenta la reflexión crítica, profundiza la comprensión del papel de la IA en el aprendizaje y ayuda a evaluar su impacto en la calidad e integridad del trabajo. Esta página puede eliminarse configurando la opción \texttt{aiack} como \texttt{false} en las opciones de clase del documento.}
    
        \vspace{2em}
        
        \exampleinfo{%
        Reconozco el uso de Gemini (\url{https://gemini.google.com}) para perfeccionar el tono académico y mejorar la precisión lingüística de este trabajo, incluyendo aspectos de gramática, puntuación y vocabulario.
        
        \vspace{1em}
        
        \noindent\textbf{Descripción del Uso}
        
        \prompt[1]{En un párrafo, describa simplemente qué es LaTeX.}
        
        \aioutput[1]{LaTeX es un sistema de composición tipográfica de alta calidad, comúnmente utilizado para producir documentos científicos y matemáticos debido a su potente manejo de fórmulas y bibliografías. A diferencia de los procesadores de texto, utiliza archivos de texto plano con etiquetas de marcado para definir la estructura y el formato de un documento, permitiendo un control preciso sobre la disposición y apariencia. Es especialmente popular en el ámbito académico y de investigación para la creación de tesis, artículos de revistas y presentaciones.}
        }
    }{%
        \chapter*{AI Acknowledgement}

        \guideinfo{It is good academic practice to briefly state how, why, and when AI was used. Explaining both the actions taken and the reasons behind them encourages critical reflection, deepens understanding of AI's role in learning, and supports the evaluation of its impact on your work's quality and integrity. This page can be removed by setting the option \texttt{aiack} to \texttt{false} in the document class options.}

        \vspace{2em}
        
        \exampleinfo{%
        I acknowledge the use of Gemini (\url{https://gemini.google.com}) to refine the academic tone and improve the linguistic accuracy of this work, including aspects of grammar, punctuation, and vocabulary.
        
        \vspace{1em}
        
        \noindent\textbf{Description of Use}
        
        \prompt[1]{In a paragraph, simply describe what LaTeX is.}
        
        \aioutput[1]{LaTeX is a high-quality typesetting system commonly used for producing scientific and mathematical documents due to its powerful handling of formulas and bibliographies. Unlike word processors, it uses plain text files with markup tags to define the structure and formatting of a document, allowing precise control over layout and appearance. It is especially popular in academia and research for creating theses, journal articles, and presentations.}
        }
    }}
    
    \MediaOptionLogicBlank   
}{}
